\textbf{Why ``verify, don't decode'' works, and where it stops.} Steganography hides the \emph{message}
but not its \emph{purpose}. Colluders coordinate in order to change an outcome: a weaker test suite, a
higher price. A monitor with independent access to that outcome is largely indifferent to how the
message was carried. This matches the information-theoretic picture of \citet{motwani2024secret}:
perfectly secure channels are undetectable \emph{as channels}, but their effects are not. The limits
follow from the same argument. (i)~If the colluders' goal leaves no trace in a verifiable outcome
(pure information leakage, as in C5b), only provenance can catch it, and provenance depends on
observability (R3). (ii)~Verification needs a ground truth the agents cannot shape. Our verifier used
upstream tests; in deployment that is a trusted QA suite, and its coverage bounds recall (R4).
(iii)~Tacit collusion (A4) is detectable as harm but cannot be attributed to an agreement. Legally and
ethically that is the right answer, and the requirements specifically ask that agreement alone never be
penalised.

\textbf{Challenges encountered.} \emph{Ground truth in real repositories:} some upstream tests import
private helpers (T3), docs-tutorial tests depend on files outside the source tree, and removing a
single hunk from a reference patch often yields code that does not import. We filtered automatically
for plausible partial fixes. \emph{History leakage:} a plain clone contains the future fix, so
workspaces are exported without history. \emph{Hard negatives create real false positives:} index
anchoring and honest price discovery look like coordination early in an episode, and online deciding
amplifies this (R5). \emph{Licensing:} the double-auction reference code has no licence, so it was
re-implemented. \emph{Tooling:} the sandbox had no third-party LLM egress or web search, which is why
M2, Petri red-teaming and LLM agents are wired in but not yet run.

\textbf{Limitations.} The central limitation is that v0.1 agents are \emph{scripted}. Their channels,
strategies and wording variation were designed by us, so the results validate the monitor's
\emph{logic and plumbing} with exact labels. They do not show that LLM colluders behave this way. In
particular, the colluding tester's knowledge of the gap is supplied by the environment as a stand-in for
a successful covert transmission. Collusion is induced, not emergent. We ran five tasks from one
repository, one market structure and three to ten seeds, so the CIs are wide and some effects are near
ceiling. Market buyers are scripted and costs are stylised. The monitor's acceptance tests derive from
upstream tests.

\textbf{Next steps.} (1)~Run the implemented LLM path: \texttt{cloud} (Claude Haiku agents, Sonnet
judge), \texttt{local} (Ollama open weights) and \texttt{mixed} teams. Test whether shared-model
Schelling points make covert schemes easier for homogeneous teams. (2)~Add \Mtwo{} and measure the
transcript judge against the same channels at a fixed FPR. (3)~Petri/MRT-style red-teaming to
generate \emph{new} schemes against the current detectors, used as held-out tests. (4)~Open-network
treatment with retrieval instrumentation, and sampling-seed channels \citep{machtle2026seeds}, which
require logging sampler seeds. (5)~Emergent collusion through retained interaction history
\citep{shi2026longhorizon}. (6)~Activation probes for open-weight agents (NARCBench-style) as an extra
provenance signal. (7)~Train a single cross-domain monitor and measure transfer.

\textbf{Conclusion.} In both a real software team and a market, a trusted monitor that independently
verifies outcomes and tracks provenance caught every coordination channel we tested at near-zero
benign cost, while a transcript monitor caught only overt collusion. Detection was channel-invariant;
attribution was not. And the best intervention was to \emph{add} a trusted participant rather than to
exclude suspects.
